\subsection*{Operations on I-Tuples}

Operations on \emph{i-tuple}s are similar to those on \emph{mdspan}. You can either use the \emph{element-of} operation \texttt{()} to retrieve the element values or use \emph{ituple} as a whole:

\begin{lstlisting}[language=croqtile,caption={Assigning a new tuple}]
a = {3, 4};
b = {a(0), 1, a(1)};  // '()' to retrieve the element value
c = a {(0), (1), 2};  // syntax sugar
d = {a, 5, 6};        // concatenate
e = a + 1;            // addition is applied elementwise
\end{lstlisting}

Similar to \emph{integer} and \emph{mdspan}, an \emph{ituple} variable must be initialized and cannot be reassigned.

In practice, \emph{ituple}s are often used with \emph{mdspan}s. For example:

\begin{lstlisting}[language=croqtile,caption={Indexing an mdspan with an i-tuple}]
shape : [7, 18, 28];
tiling_factors = {1, 2, 4};
tiled_shape : shape / tiling_factors;
padded_shape : shape + {2, 0, 2};
\end{lstlisting}

Here, dividing an \texttt{mdspan} shape by an \texttt{ituple} of tiling factors produces \texttt{tiled\_shape}. Adding an anonymous \texttt{ituple} to the original shape produces \texttt{padded\_shape}. The shape and tuple must have matching ranks; otherwise, the compiler reports an error.

\begin{lstlisting}[language=croqtile,caption={An invalid shape expression}]
shape : [7, 8, 9] + {1, 2}; // error: inconsistent rank
\end{lstlisting}
